\begin{graphillustration}{Две плоские укладки одного $K_4$: внешними стали разные грани. Метки вершин сохраняют соответствие рёбер.}{fig:14-outer-face}
\begin{tikzpicture}[font=\small]
\begin{scope}[shift={(0,0)}]
\node[gi vertex] (a) at (0,2) {$a$};
\node[gi vertex] (b) at (-1.7,-1) {$b$};
\node[gi vertex] (c) at (1.7,-1) {$c$};
\node[gi vertex] (d) at (0,0) {$d$};
\draw[gi edge] (a) -- (b);
\draw[gi edge] (a) -- (c);
\draw[gi edge] (a) -- (d);
\draw[gi edge] (b) -- (c);
\draw[gi edge] (b) -- (d);
\draw[gi edge] (c) -- (d);
\node[gi note] at (0,-1.8) {Внешняя грань: $abc$};
\end{scope}
\begin{scope}[shift={(6.3,0)}]
\node[gi vertex] (a) at (0,2) {$a$};
\node[gi vertex] (b) at (-1.7,-1) {$b$};
\node[gi vertex] (d) at (1.7,-1) {$d$};
\node[gi vertex] (c) at (0,0) {$c$};
\draw[gi edge] (a) -- (b);
\draw[gi edge] (a) -- (d);
\draw[gi edge] (a) -- (c);
\draw[gi edge] (b) -- (d);
\draw[gi edge] (b) -- (c);
\draw[gi edge] (d) -- (c);
\node[gi note] at (0,-1.8) {Внешняя грань: $abd$};
\end{scope}
\end{tikzpicture}
\end{graphillustration}
